\chapter{Related Work}
\label{chap:related_work}
In this chapter, I will go over existing approaches to visual animal re-identification (Re-ID) with an emphasis on methods intended to generalize across species. The literature will be organized into thematic categories, from traditional computer vision techniques to modern deep learning methods, including hybrid methods that combine the two. I'll also discuss key challenges in the field, such as limited and/or low quality data, pose variation and open-set identification. I'll synthesize the findings from prior work and emphasize the gaps that this thesis is trying to address.

\section{Traditional Computer Vision Approaches}
Early animal Re-ID systems relied on hand-crafted visual features and local descriptors. A prime example is HotSpotter~\cite{6475023}, a SIFT-based~\cite{Lowe2004SIFT} algorithm that identifies individuals by comparing local keypoints in images. HotSpotter uses viewpoint-invariant SIFT descriptors and scores matches by emphasizing the most distinctive keypoints (the titular "hot spots") on an animal's coat pattern. ~\cite{jmse11050889} compared SIFT, SURF and SuperPoint on identifying sunfish individuals and demonstrated that the number of matching keypoints between a pair of images is a strong indicator for the likelihood that they contain the same individual. These traditional local-feature methods do not require training on identity labels and are, by definition, species-agnostic, making them applicable to a wide variety of species. HotSpotter and similar algorithms have been successfully used on numerous species with distinctive markings, including zebras, giraffes, jaguars, ocelots and leopards.~\cite{springerSpeciesAgnosticPatterned} The drawback, however, is that purely hand-crafted methods cannot fully exploit any available training data since their matching heuristic is fixed and may miss more subtle or global cues. 

\section{Species-Specific Trait-Based Re-Identification}
Another line of research focuses on species-specific visual traits for re-identification. Instead of analyzing the entire image or generic local patches, these approaches isolate particular biometric features unique to the species at hand, for example the ear shapes of elephants, the outline of a shark's dorsal fin, or the contour of a whale's tail (fluke), and use them for re-identification. Both traditional image analysis and deep learning have been employed to recognize such traits. For example, algorithms such as CurvRank~\cite{ani12151980} and finFindR~\cite{finFindR} compute and match the shapes of whale flukes or shark fins to identify individuals. These methods often don't require individual identity labels for training, rather a detector or feature extractor for the trait can be trained using only that trait's annotations, then match those features between images. The advantage of trait-specific methods is that they leverage highly distinctive features (when they are present) and can work with minimal ID labeling. However, a major limitation is that they are species-specific so each system is tailored for identifying only one species and usually it can't be transferred to other species. This lack of generality is a significant barrier to wider use as it requires developing a new method for each studied species. This thesis is motivated by overcoming this limitation by developing a method, based on existing research, that does not depend on any single species' unique traits. 

\section{Deep Learning-Based Methods}
With the rise of deep learning, a number of methods have been proposed that learn features directly from image data, without manual feature engineering. Modern deep learning approaches for animal re-ID typically employ CNNs (Convolutional Neural Networks) or related architectures to produce an embedding (feature vector) for each image, such that images of the same individual are close in feature space while those of different individuals are far apart. These can be trained either as classification (training a CNN to classify a fixed set of known individuals) or via metric learning (training the network to differentiate between same-individual vs. different individual image pairs). In general, metric learning approaches are favored for wildlife re-ID because they naturally handle open-set problems. A classification-based model treats re-ID as identifying one of $N$ known individuals, which fails when a new individual outside the training set appears. By contrast, metric-based methods learn a similarity function and can thus match new individuals by embedding comparison, without needing to retrain for new identities.~\cite{springerSpeciesAgnosticPatterned} Recent progress in this area has been accelerated by large-scale multidataset training and more standardized benchmarking. The WildlifeDatasets toolkit~\cite{čermák2023wildlifedatasetsopensourcetoolkitanimal} provides an easy and unified access to a wide array of re-ID datasets, enabling reproducible evaluation across species. MegaDescriptor is a model trained on these datasets that out-of-the-box provides a strong backbone for re-ID tasks. At the same time, large-scale training introduced an evaluation pitfall that is not discussed enough in the literature and that is excessive data leakage due to appearance of near-duplicates in the datasets. 

\section{Hybrid Approaches with Learned Local Features}
A promising direction in animal re-id is to combine the strengths of traditional and deep learning methods into a hybrid pipeline. The idea is to use modern learned features from deep networks in place of the hand-crafted ones, but still employ them in a local matching and aggregation framework as in traditional approaches. A good example of this is the ALFRE-ID~\cite{springerSpeciesAgnosticPatterned} pipeline (Aggregated Local Features for Re-Identification). ALFRE-ID uses a CNN extractor to extract local patch descriptors from an animal's coat pattern, then aggregates these local features into a global descriptor for the image. 

WildFusion~\cite{cermak2024wildfusionindividualanimalidentification} takes a different approach by combining global deep similarity and local matching similarity using calibration to make these scores comparable. Their results have shown that calibrated fusion can substantially improve accuracy across diverse benchmarks. The trade-off is computational cost, which motivated this thesis to create a faster system while preserving accuracy. 

\section{Population Estimation}
Individual identification is something that lends itself naturally to the task of estimating animal populations. Perez et al~\cite{perez2024humanintheloopvisualreidpopulation} propose Human-in-the-Loop Importance Sampling (HITL-NIS), an estimator grounded in statistics that queries a human annotator for a relatively small number of same/different pair judgements to estimate the number of individuals without fully clustering the dataset. However, human annotation is never perfect and Johansson et al.~\cite{Johansson_Samelius_Wikberg_Chapron_Mishra_Low_2020} show that identification errors in camera trap workflows can drastically bias population estimates, emphasizing the need to account for human mistakes. This thesis directly builds on this gap by evaluating HITL-NIS under simulated annotation noise to judge the robustness of the system. 



\section{Key Challenges and Limitations}
Research in animal re-identification encounters several fundamental challenges that limit real-world performance. Key issues identified in the literature are: 

  \textbf{Limited labeled data.} Collecting individual-level labels is
   labor-intensive, so datasets often contain only a handful of images per
   subject, leading to overfitting and poor generalization.  

   \textbf{Pose variation.} Animals appear under diverse viewpoints and
   behaviors, drastically altering their visual features. Identification
   accuracy drops when profiles or unusual poses replace frontal views.  

   \textbf{Occlusions.} Vegetation, group interactions, or partial camera trap
   captures often conceal distinguishing markings, hindering feature extraction.

   \textbf{Open-set identification.} Field deployments must handle individuals
   unseen during training; closed-set models misclassify such cases. Open-set
   recognition methods highlight the need to reject unknown identities.

